% !TEX root = ../main.tex
%%%%%%%%%%%%%%%%%%%%%%%%%%%%%%%%%%%%%%%%%%%%%%%%%%%%%%%%%%%%%%%%%%%%%%%
% Discussion
%%%%%%%%%%%%%%%%%%%%%%%%%%%%%%%%%%%%%%%%%%%%%%%%%%%%%%%%%%%%%%%%%%%%%%%

\section{Discussion}
\label{sec:discussion}

\subsection{One channel, three falsified hypotheses}

The paper set out to measure a sign frontier and, in the process, found
something more robust than the frontier itself. Three separate stabilization
hypotheses were posed and each was falsified:

slower demand expectations would damp the collapse
(\cref{sec:expectations}), but they did not; a balanced-budget demand floor
would shrink the collapse region (\cref{sec:government}), but it enlarged
it, with the fiscal instrument saturated at its cap; and that same benefit
would cushion firm heterogeneity (\cref{sec:heterogeneity}), but it
lowered the viability threshold.

All three fail through the same channel and for the same reason. The
instability originates in the wage curve: unemployment moves the wage, the
wage moves profit-maximizing employment (increasingly so as $\sigma$ rises),
employment moves unemployment, and each turn of the cycle leaves investment
short of depreciation, so capital erodes. Demand-side instruments act on the
wrong variable. Worse, in case (3) the benefit fails through a mechanism that
is fully traceable and slightly perverse: by succeeding at lowering
unemployment, it causes the wage curve to raise the wage, which pushes the
weakest firm below its replacement investment.

There is a deeper reason the demand instruments miss, and it also says what would
not: the loop turns because, at a fixed numéraire, the nominal wage is the
real wage, so a wage cut is a real wage cut and no demand instrument can reach the
variable that moves. A price convention can. Relaxing the numéraire to complete
instantaneous pass-through switches the feedback off through a Pigou effect on
real balances (\cref{tab:prices}), which is why the result is best read as
bracketed by the price convention rather than settled. \Cref{sec:frontierlocal}
places that bracket alongside two further senses in which the frontier turns out
to be a local object.

The unifying boundary is worth stating plainly, because it is the most
transferable result in the paper: demand instruments work where the
binding constraint is demand, and not against wage-driven capital erosion.

\subsection{What the global analysis does and does not overturn}

The sensitivity analysis is destructive of the headline result of
\cref{sec:shape} and should be read as such: over the defensible parameter space, wage-led is the
exception; $\sigma$ is nearly inert marginally; and viability is governed by the
depreciation rate, a parameter the calibration treats as a convention rather
than an estimate (\cref{sec:calibration}). Any
version of this work that reported $\sigstar = 0.654$ as the result
would have been reporting a cell as if it were a space.

What it does not overturn is the conditional finding itself. A conditional
frontier measured at defaults and a marginal effect averaged over fifteen
other parameters are different objects, and \cref{tab:sobol} gives positive
evidence that they should differ here. The honest summary is that the model
has a robust local sign structure and a nearly inert global one.

\subsection{The measurement instrument was part of the result}

\Cref{sec:estimator} reports the uncomfortable half of the reconciliation: a
material part of the gap between the conditional and marginal readings was our
own quantity of interest. A two-point chord and a whole-support slope are
different functionals of a curved response, and on the same committed data
they place $\sigstar$ at $0.447$ and $0.643$ with non-overlapping intervals.
Re-running the global analysis on the repaired quantity moves
$P(\text{wage-led} \mid \text{viable})$ from $0.095$ to $0.026$, and the
whole of that movement is traceable to $108$ individually identified design
points, those where the turning point of the U falls inside the chord's
window.

Two lessons generalize beyond this model. First, a summary statistic
chosen before the shape of the response is known can encode a conclusion: the
chord was fixed when $Y(\rho)$ was assumed monotone, and it kept returning
answers about a curve it could not represent. Second, the fix was not to pick
the better statistic but to stop compressing a curve into one number.
That is why \cref{sec:shape} leads with the turning point and the side the
anchored range falls on, and treats the frontier as a declared reduction of
it: of the three summaries, only the position relative to the turn is
invariant to the analyst's choice of interval.

The episode has a plainer companion, recorded in \cref{app:retractions}: an
``inverted direction'' between the conditional and marginal analyses, which
this project reported as an open contradiction, did not exist in any
measurement; it was an error in our own prose. The experiment designed to
explain it was still worth running, because it found a real defect on its other
axis. But the pairing is the lesson: a project's summary can contradict its own
tables, and nothing except re-deriving each claim from the measurements will
catch it.

\subsection{The sign frontier is local; the anchored margin is not}
\label{sec:frontierlocal}

Three findings, on three different designs, converge on the same qualification of
the sign frontier, and are worth stating as one. Each shows the frontier moving,
or dissolving, when a background choice is varied; none of them touches the result
of \cref{sec:shape}.

First, the frontier is conditional on fiscal institutions: the balanced-budget
benefit of \cref{sec:government} eliminates the wage-led region entirely at an
OECD-band replacement rate. Second, it is nearly inert marginally: with the other
parameters swept, the first-order sensitivity of the whole-support slope to
$\sigma$, conditioned on viability, is $0.019$ (\cref{sec:sa}), a frontier
defined by a sign crossing in $\sigma$ barely responds to $\sigma$ once the rest
of the space is allowed to move. Third, it is conditional on the price numéraire:
the pass-through probe of \cref{sec:wagecurve} lifts $\sigstar$ upward, by roughly
a grid node in $\sigma$ at moderate wage flexibility and more at high $\eta$
(\cref{tab:prices}).

What survives all three is not the position of the frontier but the
anchored margin of the first result. Across fiscal settings, across the
marginalized space, and across both price conventions, the anchored retention
rate stays to the left of the turning point for $\sigma \ge 0.4$, so the
marginal effect of retention at the anchored point keeps its sign. The frontier
is a local, conditional object; the side the anchored range falls on is the robust
one. Read this way the paper has a hierarchy of its own results, a fragile
headline ($\sigstar$ as a number) resting on a durable one (the anchored margin of
the U-shape), rather than a single frontier asked to carry more than a
conditional cell can.
